\begin{lemma}[The hard truncation is redundant when the input vanishes near the ball]
Fix $n \geq 1$ and $j \in \{1,\ldots,n\}$, let $\lambda$ be Lebesgue measure on
$\mathbb{R}^n$, let $\gamma_n$ be the Gaussian measure of \noderef{GaussMeasure}, and let
$K_{\gamma,j}$ be the normalized Gaussian Riesz kernel of \noderef{GaussRieszKernel}.  Fix
$x_0 \in \mathbb{R}^n$ and $\rho > 0$, and for $z \in \mathbb{R}^n$ write
\[
  \Omega_z = \{y \in \mathbb{R}^n : |z-y| > \rho/4\} .
\]
Let $h : \mathbb{R}^n \to \mathbb{R}$ be $\gamma_n$-integrable and suppose
\[
  h = 0 \quad \lambda\text{-almost everywhere on } B(x_0,\tfrac{3\rho}{4}) .
\]
Then for \emph{every} $z \in B(x_0,\rho/2)$ the function $y \mapsto K_{\gamma,j}(z,y)h(y)$ is
Lebesgue integrable on $\mathbb{R}^n$, and
\[
  \int_{\mathbb{R}^n}K_{\gamma,j}(z,y)\,h(y)\,dy
    = \int_{\Omega_z}K_{\gamma,j}(z,y)\,h(y)\,dy .
\]

\emph{Both conclusions are content, and the first is not an afterthought.}  The two displayed
integrals are Bochner integrals, each equal to $0$ by convention on a non-integrable integrand;
an equality between them could therefore hold with both sides conventional and assert nothing
about the kernel.  The integrability assertion is what rules that out, and it is stated for
\emph{every} $z$ in the ball rather than almost every one, so that a consumer may evaluate
either side pointwise.

\emph{The two radii are not interchangeable.}  The hypothesis is vanishing on the ball of
radius $3\rho/4$ and the conclusion holds for centres in the ball of radius $\rho/2$, with the
truncation taken at $\rho/4$; the three are tied by $\rho/4 + \rho/2 = 3\rho/4$, which is
exactly what makes the discarded region lie inside the region where $h$ vanishes.  Weakening the
hypothesis to vanishing on $B(x_0,\rho/2)$, or widening the conclusion to centres in
$B(x_0,3\rho/4)$, would break the proof and the statement.

The statement is deliberately about an arbitrary such $h$ rather than about one particular
function: its consumer applies it both to a given $L^1(\gamma_n)$ input and to each member of a
sequence of smooth approximants to that input, and those are different functions sharing only
the vanishing hypothesis.
\end{lemma}
\begin{proof}
Put $R = \|x_0\| + \rho + 1 > 0$, and let $C > 0$ be the constant supplied by clause (1) of
\noderef{GaussRieszKernelFarMajorant} for the pair $(R,\rho/4)$, both of which are positive.
Fix $z \in B(x_0,\rho/2)$, so that $|z - x_0| < \rho/2$.

\emph{Step 1: $z$ lies in the ball the far majorant needs.}
By the triangle inequality, $|z| = |(z-x_0) + x_0| \leq |z-x_0| + |x_0| < \rho/2 + \|x_0\|$,
and $\rho/2 + \|x_0\| \leq \|x_0\| + \rho + 1 = R$ because $\rho/2 \leq \rho + 1$.  So
$|z| \leq R$.

\emph{Step 2: integrability on $\Omega_z$.}
The set $\Omega_z = \{y : \rho/4 < |z-y|\}$ is measurable, being the preimage of an open
half-line under the continuous map $y \mapsto |z-y|$.  Clause (1)(b) of
\noderef{GaussRieszKernelFarMajorant}, applied to the $\gamma_n$-integrable function $h$ and to
this $z$ --- legitimate by Step 1 --- states that $y \mapsto K_{\gamma,j}(z,y)h(y)$ is Lebesgue
integrable on $\{y : \rho/4 < |z-y|\} = \Omega_z$.  (That clause also supplies a bound on the
integral, which is not needed here.)

\emph{Step 3: the integrand vanishes $\lambda$-almost everywhere off $\Omega_z$.}
Let $y \notin \Omega_z$, that is $|z-y| \leq \rho/4$.  Then, using $|y-z| = |z-y|$,
\[
  |y-x_0| = |(y-z) + (z-x_0)| \leq |y-z| + |z-x_0| < \tfrac{\rho}{4} + \tfrac{\rho}{2}
    = \tfrac{3\rho}{4} ,
\]
so $y \in B(x_0,3\rho/4)$.  By hypothesis $h$ vanishes $\lambda$-almost everywhere on that
ball; hence for $\lambda$-almost every $y$, if $y \notin \Omega_z$ then $h(y) = 0$ and therefore
$K_{\gamma,j}(z,y)h(y) = 0$.

\emph{Step 4: conclusion.}
Write $F(y) = K_{\gamma,j}(z,y)h(y)$ and let $\mathbf 1_{\Omega_z}F$ be the function equal to
$F$ on $\Omega_z$ and to $0$ elsewhere.  Since $\Omega_z$ is measurable, $\mathbf 1_{\Omega_z}F$
is Lebesgue integrable on $\mathbb{R}^n$ if and only if $F$ is Lebesgue integrable on
$\Omega_z$ (shared context: the indicator criterion for integrability on a measurable set); by
Step 2 it therefore is.  By Step 3, $F = \mathbf 1_{\Omega_z}F$ $\lambda$-almost everywhere: the
two agree at every point of $\Omega_z$ by definition, and at almost every point off $\Omega_z$
both vanish.  Two functions agreeing almost everywhere are integrable together, so $F$ is
Lebesgue integrable on $\mathbb{R}^n$.  This is the first conclusion.

For the second, Step 3 says precisely that $F$ vanishes almost everywhere on the complement of
$\Omega_z$, and for such an $F$ the integral over $\Omega_z$ equals the integral over the whole
space (shared context; this identity needs no integrability hypothesis, both sides being read
with the convention, and it is applied here to an $F$ already known integrable).  That is the
displayed equality, read from right to left.
\end{proof}
